\chapter{Implementaci\'on}
El proyecto se ha desarrollado siguiendo una arquitectura modular que permite separar claramente las distintas funcionalidades relacionadas con el procesamiento de im\'agenes mediante aut\'omatas celulares. A continuaci\'on, se detallan los principales componentes te\'oricos e implementativos y c\'omo se integran en el sistema:

\section{Fundamentos Te\'oricos}
\subsection{Aut\'omatas Celulares y Regla 90R}
El sistema implementa una variante de la Regla 90 de Wolfram, denominada Regla 90R (Reversible), que constituye un aut\'omata celular bidimensional con propiedades criptogr\'aficas. La regla opera sobre un arreglo binario donde cada celda evoluciona seg\'un la funci\'on:
\begin{equation}
C_{i,j} = (C_{i-1,j-1} + C_{i-1,j} + C_{i-1,j+1}) \bmod 2
\end{equation}
donde $C_{i,j}$ representa el estado futuro de la celda en la posici\'on $(i,j)$, y la suma se realiza m\'odulo 2 (operaci\'on XOR). La implementaci\'on utiliza \'indices alternos para garantizar la reversibilidad del proceso.

\subsection{Representaci\'on Binaria de Im\'agenes}
Las im\'agenes en escala de grises se transforman en arreglos binarios mediante el desempaquetado de bits (\texttt{np.unpackbits}), donde cada p\'ixel de 8 bits se convierte en 8 celdas binarias independientes. Esta representaci\'on permite la aplicaci\'on directa de operaciones de aut\'omatas celulares sobre los datos de imagen.

\section{M\'odulos del Sistema}
\subsection{Conversi\'on de Imagen}
\textbf{Funciones}: \texttt{image\_to\_bit\_array}, \texttt{bit\_array\_to\_image}\\
\textbf{Ubicaci\'on}: Regi\'on de funciones de conversi\'on\\
\textbf{Descripci\'on}:

La funci\'on \texttt{image\_to\_bit\_array} convierte una imagen en escala de grises a una representaci\'on binaria utilizando el siguiente procedimiento:
\begin{enumerate}
\item Carga la imagen usando \texttt{PIL.Image.open()} y la convierte a escala de grises con \texttt{convert("L")}.
\item La imagen se convierte a un arreglo NumPy mediante \texttt{np.array()}.
\item El arreglo de bytes se desempaqueta en bits con \texttt{np.unpackbits()} a lo largo del eje de columnas.
\item Se devuelve el arreglo binario junto con los metadatos necesarios (ancho, alto, padding) para su reconstrucci\'on.
\end{enumerate}

La funci\'on inversa \texttt{bit\_array\_to\_image} reconstruye la imagen original a partir del arreglo binario:
\begin{enumerate}
\item El arreglo se recorta para eliminar el padding y coincidir con el ancho original.
\item Se empaquetan los bits nuevamente con \texttt{np.packbits()}.
\item El arreglo se reinterpreta como una imagen y se convierte a objeto \texttt{PIL.Image} en modo \texttt{L} (escala de grises).
\end{enumerate}

\subsection{L\'ogica del Aut\'omata Celular}
\subsubsection{Codificaci\'on}
\textbf{Funci\'on}: \texttt{encode\_rule\_90R}\\
\textbf{Descripci\'on}:

Esta funci\'on implementa la evoluci\'on del aut\'omata seg\'un la Regla 90R en modo codificaci\'on. Se trabaja sobre las filas impares del arreglo binario, desplazando los valores de las filas adyacentes para aplicar la operaci\'on XOR de vecindarios. Las operaciones se vectorizan con NumPy para mayor eficiencia.

\begin{algorithm}[H]
\caption{Algoritmo de Codificaci\'on Regla 90R}
\begin{algorithmic}[1]
\FOR{$iter = 1$ \TO $iterations$}
\STATE $new\_bits \gets bits.copy()$
\STATE $i\_idx \gets \{1, 3, 5, \ldots, rows-1\}$ \COMMENT{\'Indices impares}
\STATE Calcular vecindarios izquierdo y derecho
\FOR{$i \in i\_idx$}
\STATE $new\_bits[i, :] \gets (left[i, :] + bits[i-1, :] + right[i, :]) \bmod 2$
\STATE $new\_bits[i-1, :] \gets bits[i, :]$
\ENDFOR
\STATE $bits \gets new\_bits$
\ENDFOR
\end{algorithmic}
\end{algorithm}

\subsubsection{Decodificaci\'on}
\textbf{Funci\'on}: \texttt{decode\_rule\_90R}\\
\textbf{Descripci\'on}:

Este procedimiento invierte el proceso de codificaci\'on. Se opera ahora sobre las filas pares, usando el estado de las filas inferiores y vecinos laterales. La decodificaci\'on es sim\'etrica al proceso anterior, pero en direcci\'on inversa.

\begin{algorithm}[H]
\caption{Algoritmo de Decodificaci\'on Regla 90R}
\begin{algorithmic}[1]
\FOR{$iter = 1$ \TO $iterations$}
\STATE $new\_bits \gets bits.copy()$
\STATE $i\_idx \gets \{0, 2, 4, \ldots, rows-2\}$ \COMMENT{\'Indices pares}
\STATE Calcular vecindarios izquierdo y derecho
\FOR{$i \in i\_idx$}
\STATE $new\_bits[i, :] \gets (left[i, :] + bits[i+1, :] + right[i, :]) \bmod 2$
\STATE $new\_bits[i+1, :] \gets bits[i, :]$
\ENDFOR
\STATE $bits \gets new\_bits$
\ENDFOR
\end{algorithmic}
\end{algorithm}

\subsection{Interfaz Gr\'afica de Usuario}
\subsubsection{Arquitectura de la GUI}
\textbf{Clase}: \texttt{AutomataGUI}\\
\textbf{Framework}: Tkinter con TTK\\
\textbf{Descripci\'on}:

La interfaz gr\'afica encapsula toda la l\'ogica del sistema interactivo en una clase basada en Tkinter. Ofrece:
\begin{itemize}
\item Control de par\'ametros mediante \texttt{tk.IntVar}.
\item Selecci\'on de archivos con \texttt{filedialog.askopenfilename}.
\item Visualizaci\'on de resultados con \texttt{matplotlib.pyplot}.
\item Manejadores de errores mediante \texttt{try-except} y \texttt{messagebox}.
\end{itemize}

\subsubsection{Flujo de Operaci\'on del Emisor}
\textbf{M\'etodo}: \texttt{run\_emisor}

Este m\'etodo ejecuta:
\begin{enumerate}
\item Carga de la imagen de entrada.
\item Conversi\'on a bits y codificaci\'on con la regla 90R.
\item Almacenamiento del resultado codificado como archivo \texttt{.npy}.
\item Visualizaci\'on comparativa de imagen original y cifrada.
\end{enumerate}

\subsubsection{Flujo de Operaci\'on del Receptor}
\textbf{M\'etodo}: \texttt{run\_receptor}

Este m\'etodo invierte el proceso:
\begin{enumerate}
\item Carga del archivo \texttt{.npy} con los datos codificados.
\item Extracci\'on de metadatos.
\item Aplicaci\'on del proceso de decodificaci\'on.
\item Reconstrucci\'on de la imagen y visualizaci\'on del resultado.
\end{enumerate}

\section{Visualizaci\'on y An\'alisis}
\subsection{Visualizaci\'on de Resultados}
\textbf{Funci\'on}: \texttt{visualize\_images}

Esta funci\'on utiliza Matplotlib para mostrar comparaciones entre im\'agenes, configurando los subplots din\'amicamente seg\'un el n\'umero de im\'agenes. Elimina ejes para claridad y agrega t\'itulos descriptivos.

\subsection{Almacenamiento de Datos}
Los datos binarios intermedios se almacenan en formato \texttt{.npy}, lo cual permite:
\begin{itemize}
\item Eficiencia de lectura/escritura.
\item Preservaci\'on de dimensiones originales.
\item Alta compatibilidad con herramientas de procesamiento en Python.
\end{itemize}

\section{Consideraciones de Implementaci\'on}
\subsection{Optimizaci\'on Computacional}
El c\'odigo se ha optimizado mediante:
\begin{itemize}
\item Vectorizaci\'on con NumPy.
\item Preasignaci\'on de memoria para arreglos.
\item Reducci\'on de copias innecesarias mediante operaciones \texttt{in-place}.
\end{itemize}

\subsection{Manejo de Errores}
El sistema incluye validaci\'on exhaustiva:
\begin{itemize}
\item Verificaci\'on de tipo y formato de entradas.
\item Comprobaci\'on de consistencia antes/despu\'es del cifrado.
\item Notificaciones al usuario mediante ventanas emergentes.
\end{itemize}

\subsection{Escalabilidad}
La estructura modular admite:
\begin{itemize}
\item Inclusi\'on de nuevas reglas criptogr\'aficas.
\item Reemplazo f\'acil de componentes.
\item Reutilizaci\'on del c\'odigo base en otros proyectos.
\end{itemize}
